\documentclass[pdflatex,sn-mathphys-num]{sn-jnl}

\usepackage{graphicx}
\usepackage{multirow}
\usepackage{amsmath,amssymb,amsfonts}
\usepackage{amsthm}
\usepackage{mathrsfs}
\usepackage[title]{appendix}
\usepackage{xcolor}
\usepackage{textcomp}
\usepackage{manyfoot}
\usepackage{booktabs}
\usepackage{algorithm}
\usepackage{algorithmicx}
\usepackage{algpseudocode}
\usepackage{listings}

\raggedbottom

\begin{document}

\title[Master's Thesis Proposal]{Beyond Pairwise GUI Quality Judgments: Multi-Dimensional Evaluation with Modern Multimodal LLMs}

\author{\fnm{Xinyang} \sur{Xie}}

\affil{\orgdiv{Department of Informatics}, \orgname{Technische Universitaet Clausthal},
\orgaddress{\city{Clausthal-Zellerfeld}, \country{Germany}}}

\abstract{
Assessing the quality of graphical user interfaces (GUIs) is becoming increasingly
important as large language models are used to generate interface designs and frontend
code. Existing screenshot-based approaches can support GUI quality judgments, but they
remain limited by binary decisions, older models, and simple prompting strategies.
This proposal extends that line of work with current multimodal large language models.
The evaluation is organized around two complementary dataset tracks: a
screenshot-based baseline drawn from existing UI resources, and a requirement-driven
track in which interfaces are generated from natural language descriptions. The study
combines pairwise comparison with rubric-based assessment across multiple dimensions,
including layout quality, visual hierarchy, information organization, perceived
usability, and -- for generated interfaces -- requirement fidelity and interaction
quality. Methodologically, the work begins with a literature review of
LLM-as-a-judge techniques to identify applicable strategies, followed by a systematic
comparison of approaches such as zero-shot prompting, rubric-guided scoring, repeated
judging, and order-swapped pairwise comparisons, where the order of candidates is
reversed to detect positional bias. A human annotation protocol with agreement
analysis provides reference labels. The resulting protocol will be packaged into a
benchmark and lightweight leaderboard that can be extended to newly added models. In
addition to the static evaluation tracks, the proposal includes a small dynamic
evaluation module in which computer-use agents attempt short tasks on selected
interfaces. The goal is to determine which evaluation setup aligns best with human
judgments and how static quality judgments relate to dynamic interaction performance.
}

\maketitle

%%------------------------------------------------------------
\section{Introduction and Motivation}\label{sec:intro}
%%------------------------------------------------------------

The quality of a graphical user interface (GUI) shapes how efficiently users complete
tasks and how they perceive a software system. As generative models are increasingly
used to produce interface mockups and frontend code, the question of how UI quality can
be evaluated becomes more pressing. Importantly, GUI quality is not limited to visual
appearance alone: it also concerns whether an interface supports successful interaction
and task completion once users begin to navigate it.

UIClip~\cite{uiclip} provided an important first step by showing that vision-language
models can perform pairwise comparisons of UI screenshots and identify which of two
interfaces appears better. This screenshot-based setting is well suited for evaluating
visual structure and perceived usability, but it does not capture whether an interface
also supports reliable task completion in interactive use. For that aspect, dynamic
assessment methods, such as computer-use agents that can interact with interfaces
automatically, are particularly relevant.

Recent work on improving UI generation from designer feedback~\cite{designer_feedback}
demonstrates that high-quality, workflow-aligned feedback can be far more informative
than large amounts of coarse ranking data; notably, simple pairwise rankings in that
study reached only around 50\% evaluator agreement, close to chance level. These
findings suggest that GUI quality assessment benefits from richer evaluation
structures -- such as rubric-based scoring and requirement-aligned judgment -- rather
than binary preference alone. In parallel, the broader LLM-as-a-judge
literature~\cite{llm_judge} has introduced prompting, consistency, and aggregation
strategies that may also be useful for GUI evaluation. Taken together, these
developments motivate a more methodologically explicit re-evaluation of GUI quality
judgment with contemporary multimodal models.

This thesis therefore asks whether modern multimodal large language models can
approximate human judgments of GUI quality beyond pairwise comparison, and how such
judgments should be elicited and evaluated in a reliable way. Its main focus is a
benchmark and leaderboard for GUI quality evaluation built around two complementary
tracks: a screenshot-based baseline track and a requirement-driven generation track.
This benchmark is complemented by a small dynamic evaluation module in which
computer-use agents examine task completion on selected interfaces.

%%------------------------------------------------------------
\section{Planned Approach}\label{sec:approach}
%%------------------------------------------------------------

The planned work centers on an evaluation pipeline with six main components: dataset
construction, human annotation, rubric design, the comparison of LLM-as-a-judge
strategies, benchmark and leaderboard design, and a small dynamic evaluation
extension.

\subsection{Datasets and Evaluation Material}

The evaluation material will be organized into two complementary tracks. The first
track draws on publicly available screenshot-based resources from UIClip as an
established baseline. These existing screenshots will be used for pairwise comparison
and rubric-based assessment of visual design quality. To avoid evaluating only older
interfaces, this material may be supplemented by a smaller curated set of contemporary
web or mobile screenshots.

The second track will be constructed around natural language requirements. A set of
short UI descriptions will be collected or designed, specifying the intended
functionality, layout, and content of an interface. For this track, a small set of
interfaces will be generated from shared natural language descriptions. Recent work
has shown that current multimodal LLMs can generate renderable front-end code from
visual and textual UI specifications~\cite{design2code}, making this pipeline
practically feasible. The resulting front-end code will then be rendered as screenshots
for visual evaluation, while remaining executable for the dynamic evaluation module.
Because this track produces runnable front-end code rather than static images alone,
it supports both browser-based interaction tests and checks of whether the generated
interface satisfies the original requirements.

Together, the two tracks allow the benchmark to cover both existing, human-designed
interfaces and newly generated ones, broadening the range of evaluation scenarios
beyond what a screenshot-only dataset would permit.

\subsection{Human Annotation Protocol}

Human judgments will be collected for a selected evaluation subset. For pairwise tasks,
annotators will compare two GUI screenshots and indicate which interface is better
overall, with an optional tie label if required. For single-image tasks, annotators
will score each GUI with a small rubric and provide brief justifications. To improve
the reliability of the reference labels, each item will be annotated by multiple
raters, and inter-rater reliability will be measured, for example with
Krippendorff's alpha.

\subsection{Multi-dimensional GUI Rubric}

To move beyond binary comparison, I will define a small rubric grounded in common
visual design principles and usability heuristics. The initial dimensions are:

\begin{itemize}
    \item \textbf{Layout quality}: spatial organization, alignment, spacing, and
    overall balance.
    \item \textbf{Visual hierarchy}: whether contrast, size, and color guide attention
    appropriately.
    \item \textbf{Information organization}: clarity of grouping and logical ordering
    of content.
    \item \textbf{Perceived usability}: how easily key actions and interface purpose
    can be inferred from a static screenshot.
    \item \textbf{Requirement fidelity} (Track B only): the degree to which the
    generated interface realizes the elements, structure, and functionality specified
    in the original natural language description.
    \item \textbf{Interaction quality} (Track B only): responsiveness to user actions,
    appropriateness of feedback and state transitions, and overall smoothness of the
    interactive experience when the generated interface is executed in a browser.
\end{itemize}

This list represents an initial set of dimensions; additional dimensions may emerge
during pilot annotation and will be incorporated if they prove reliable and
informative. Each dimension will be scored on a fixed ordinal scale with short written
anchors so that the same rubric can be applied by both human annotators and model
judges. The first four dimensions can be assessed from static screenshots alone. The
requirement fidelity dimension applies exclusively to the requirement-driven track,
where an original description is available as a reference; for this dimension, judges
compare the rendered interface against the description and assess how completely and
accurately the requirements have been met. The interaction quality dimension likewise
applies only to the requirement-driven track and will be evaluated through the dynamic
evaluation module described in Section~2.6, where computer-use agents execute short
tasks on the running interface.

\subsection{Prompting and Judging Strategies}

The selection of judging strategies will be informed by an initial literature review of
LLM-as-a-judge methods. In particular, this review will examine foundational work on
pairwise judgment and judge bias~\cite{llm_judge}, broader taxonomies of
LLM-as-a-judge techniques~\cite{li-etal-2025-generation}, and recent evidence on
position bias in pairwise evaluation~\cite{shi-etal-2025-judging}. Based on this
review, several judging strategies will be compared. A zero-shot baseline will ask the
model directly for a pairwise preference or a rubric-based score. A few-shot variant
will add one or two worked examples to test whether light calibration improves judgment
quality. A rubric-guided condition will provide explicit scoring instructions and
require a structured output format. To test robustness, the same prompt will be
repeated to measure self-consistency. For pairwise tasks, order-swapped pairwise
comparisons will be used, where the order of candidates is reversed to detect
positional bias. In addition, the study will compare single-model judgments with simple
aggregation strategies such as majority voting for pairwise tasks and averaged scores
for rubric-based tasks.

\subsection{Benchmark and Leaderboard Design}

Beyond comparing individual models, the project will package the evaluation setup into
a benchmark structure that combines pairwise tasks, rubric-based scoring, and shared
reference annotations. Results will be stored in a standardized format so that newly
added multimodal models can be evaluated under the same protocol. As a practical
artifact, the benchmark will be accompanied by a lightweight leaderboard or web-based
dashboard for comparing models across tasks and rubric dimensions.

\subsection{Dynamic Evaluation with Computer-Use Agents}

To complement the static benchmark, the project will include a small dynamic evaluation
extension. This module will operate primarily on the requirement-driven track, where
LLM-generated HTML interfaces can be rendered and executed in a browser. For a selected
subset of these interfaces, short interaction tasks derived from the original natural
language descriptions will be defined and executed by computer-use agents. Prior work
has demonstrated that multimodal agents can perform goal-directed interaction tasks on
visual web interfaces~\cite{visualwebarena}, providing a methodological foundation for
this module. This extension will be kept deliberately small and will focus on whether
interfaces that receive strong static quality judgments also support reliable task
completion in practice, and whether requirement alignment in the static setting predicts
successful interaction in the dynamic setting.

%%------------------------------------------------------------
\section{Evaluation}\label{sec:evaluation}
%%------------------------------------------------------------

Several current multimodal models will be tested as judges, including representative
models from major closed-source and open multimodal model families, depending on
availability at evaluation time. As an exploratory comparison, the study may also
include one or two smaller or weaker multimodal models in order to examine how judge
quality varies with model capability. The pairwise setting will be evaluated through
agreement with human majority labels, while the rubric-based setting will be evaluated
through rank correlation and score differences between model outputs and human ratings.
For the requirement-driven track, an additional evaluation dimension will measure how
well model judgments capture whether a generated interface satisfies its original
natural language description. The dynamic extension will be evaluated through task
completion rate, action efficiency, and qualitative failure analysis on the selected
interactive tasks.

Beyond task accuracy, the evaluation will explicitly examine reliability. This includes
inter-rater agreement in the human annotations, consistency across repeated model
judgments, stability under order-swapped pairwise comparisons, and agreement across
different models. Computational cost will also be tracked in order to assess practical
trade-offs between evaluation quality and API usage. A further goal is to determine
whether the resulting benchmark and leaderboard provide a stable and extensible basis
for comparing both current and newly introduced multimodal models.

Finally, a qualitative error analysis will examine typical failure cases, such as
cluttered layouts, minimal interfaces, text-heavy screens, and visually appealing but
less usable designs. Together, these analyses should clarify how well modern
multimodal large language models can serve as judges of GUI quality, which rubric
dimensions are most reliable, and which judging strategies align best with human
assessment.

\backmatter

{\footnotesize\setlength{\bibsep}{0pt}\bibliography{references}}

\end{document}
